\documentclass[journal]{IEEEtai}

\usepackage{amsmath}
\usepackage{amssymb}
\usepackage{amsfonts}
\usepackage{booktabs}
\usepackage{multirow}
\usepackage{url}
\usepackage[colorlinks,urlcolor=blue,linkcolor=blue,citecolor=blue]{hyperref}
\usepackage{xcolor}
\usepackage{graphicx}

\title{MoWM: Mixture of World Models via Hypernetwork-Generated VSA Transitions}

\author{Nicolas Cloutier\thanks{N. Cloutier is with the Grillcheese Research Laboratory, Levis, QC, Canada (ORCID: 0009-0008-5289-5324).}\thanks{DOI: \url{https://doi.org/10.5281/zenodo.19192788}}}

\newtheorem{theorem}{Theorem}
\newtheorem{lemma}{Lemma}

\begin{document}

\markboth{IEEE Transactions on Artificial Intelligence}{Nicolas Cloutier: MoWM: Mixture of World Models via Hypernetwork-Generated VSA Transitions}

\maketitle

\begin{abstract}
World models that learn environment dynamics as a single transition function suffer from poor generalization, compounding rollout errors, and inability to represent multi-modal dynamics. We propose MoWM (Mixture of World Models), an architecture where $M$ hypernetworks each generate a different VSA-based transition function, and a learned gating mechanism selects which world models apply per situation. Each world model operates in block-code vector space where transitions are element-wise binding operations---algebraically structured, magnitude-preserving, and invertible. The gating mechanism uses DSelect-k sparse expert selection, trained end-to-end via DisARM gradient estimation through discrete model selection. A contrastive value estimator (CVL) provides TD-free Q-values from the mixture's occupancy measure, enabling planning without explicit reward modeling. We present the MoWM architecture and validate its components individually: block-code binding at 57~$\mu$s/op with 4~$\mu$s/op similarity, FNet mixing 2$\times$ faster than Combiner-Axial attention, HYLA hypernetworks variance-stable through 5 layers, DisARM gradient variance 6.4$\times$ lower than REINFORCE, and VS-Graph achieving 100\% classification accuracy on star-vs-chain topology tasks.
\end{abstract}

\begin{IEEEImpStatement}
World models underpin autonomous planning in robotics, game AI, and decision-support systems. Current approaches rely on monolithic neural transition functions that fail in multi-modal environments and compound errors over long planning horizons. MoWM introduces algebraically structured mixtures of world models that preserve invertibility and compositionality, enabling reliable multi-step rollouts on consumer hardware drawing under 25 watts. By replacing datacenter-scale model-based reinforcement learning with sparse expert selection over block-code transitions, the technology reduces both energy consumption and hardware requirements by orders of magnitude. Applications include low-power autonomous navigation, interpretable planning for safety-critical systems, and efficient simulation for resource-constrained edge devices.
\end{IEEEImpStatement}

\begin{IEEEkeywords}
Vector symbolic architecture, world models, mixture of experts, hypernetworks, block codes, reinforcement learning, compositional planning
\end{IEEEkeywords}

\section{Introduction}

\IEEEPARstart{A}{} fundamental challenge in artificial intelligence is learning internal models of the world that support planning, counterfactual reasoning, and generalization to novel situations. Recent work on world models~\cite{ha2018world, hafner2020mastering, hansen2023tdmpc2} has achieved impressive results by learning transition functions $T: \mathcal{S} \times \mathcal{A} \rightarrow \mathcal{S}$ as neural networks, but these approaches suffer from three limitations:

\begin{enumerate}
\item \textbf{Single transition function.} Real environments have multi-modal dynamics---the same action can have different effects depending on context (e.g., pushing a door that is locked vs.\ unlocked). A single $T$ cannot naturally represent this without growing exponentially in capacity.

\item \textbf{Compounding rollout errors.} When the learned $T$ is iteratively applied for multi-step planning, small errors accumulate exponentially. Without algebraic structure, there is no natural error-correction mechanism.

\item \textbf{No compositional generalization.} Standard neural $T$ functions cannot compose actions algebraically---predicting the effect of action sequences $A_1, A_2, A_3$ requires sequential forward passes, not a single composed operation.
\end{enumerate}

Chung et al.~\cite{chung2025geometric} addressed limitations (2) and (3) by using VSA binding as the transition: $S_{t+1} = \phi(S_t) \odot \phi(A_t)$, where $\odot$ is element-wise complex multiplication in FHRR space. This provides exact invertibility ($S_t = S_{t+1} \odot \phi(A_t)^{-1}$), multi-step composition ($S_{t+k} = S_t \odot \prod \phi(A_{t+j})$), and cleanup-based error correction. However, their approach uses a single encoder pair, limiting representational capacity.

We address all three limitations with \textbf{MoWM}: a mixture of $M$ hypernetwork-generated world models, each producing a different transition function in block-code VSA space. Our key insight is that each HMM-VSA rule~\cite{grillcheese2025hmm} already \emph{is} a world model---the transition matrix $A$ encodes how states evolve under different patterns. By generating each rule's transition weights via a separate HYLA hypernetwork, we obtain $M$ specialized world models that can represent multi-modal dynamics.

\subsection{Contributions}

\begin{itemize}
\item \textbf{MoWM architecture}: $M$ HYLA hypernetworks generate $M$ different block-code transition functions. A DSelect-k gate selects the top-$k$ models per state-action pair, combining their predictions via learned mixing weights. The mixture prediction is:
\begin{equation}
S_{t+1} = \sum_{m=1}^{M} w_m(S_t, A_t) \cdot \mathrm{bind}(S_t, \Delta_m(S_t, A_t))
\end{equation}
where $\Delta_m$ is the $m$-th hypernetwork's predicted transformation and $w_m$ are DSelect-k gate weights.

\item \textbf{Automatic specialization}: Without supervision, different world models learn different aspects of the dynamics---one may specialize in spatial transitions, another in temporal patterns, a third in causal relationships. This emerges from the diversity loss on HMM transition matrices combined with DSelect-k's entropy regularizer.

\item \textbf{Algebraic planning via composed world models}: For multi-step planning, we compose the selected world models' transformations algebraically:
\begin{equation}
S_{t+k} = S_t \odot \prod_{j=1}^{k} \Delta_{m^*_j}(S_{t+j-1}, A_{t+j-1})
\end{equation}
where $m^*_j$ is the selected model at step $j$. This is $O(k)$ binding operations, not $O(k)$ full forward passes.

\item \textbf{Contrastive value estimation over the mixture}: CVL estimates Q-values from the mixture's discounted occupancy measure without TD learning. The HMM forward algorithm propagates state probabilities through the selected world models, and the occupancy measure is reward-weighted to produce Q-values.

\item \textbf{Theoretical grounding}: We connect MoWM to the GFSA framework~\cite{johnson2020learning}, showing that each world model is a finite-state policy on a graph-based POMDP, and the mixture is a hierarchical GFSA with the DSelect-k gate as the meta-policy.
\end{itemize}

\section{Background}

\subsection{VSA-Based World Models}

Chung et al.~\cite{chung2025geometric} encode states and actions as unitary complex vectors via learnable FHRR encoders:
\begin{equation}
\phi_S(s) = \left[e^{i\theta_{j,s}^\top s}\right]_{j=1}^D, \quad \phi_A(a) = \left[e^{i\theta_{j,a}^\top a}\right]_{j=1}^D
\end{equation}

Transitions are modeled as element-wise binding:
\begin{equation}
\phi_S(s_{t+1}) = \phi_S(s_t) \odot \phi_A(a_t)
\end{equation}

Training minimizes binding loss $\mathcal{L}_\mathrm{bind} = \|\phi(s') - \phi(s) \odot \phi(a)\|^2$ with invertibility and orthogonality regularizers. Cleanup via codebook similarity search corrects accumulated errors during multi-step rollouts.

\textbf{Limitation}: A single $\phi_A$ encoder maps each action to a fixed transformation, independent of context. In environments where the same action has different effects depending on state (e.g., walls, locked doors, varying terrain), this fixed mapping cannot adapt.

\subsection{Block Codes vs FHRR}

Chung et al.\ use FHRR (continuous complex phase vectors). We use sparse block codes~\cite{hersche2023nvsa, grillcheese2025hmm} where each vector consists of $k$ blocks of length $l$, with one-hot encoding per block. Block codes provide:
\begin{itemize}
\item \textbf{Exact magnitude preservation} under chained binding (FHRR drifts).
\item \textbf{Natural hashing} for HyperAttention (block argmax = SimHash bucket).
\item \textbf{Discrete cleanup} via per-block argmax (faster than continuous similarity search).
\end{itemize}

\subsection{HYLA Hypernetworks}

HYLA~\cite{grillcheese2025hmm} generates weight matrices from block-code conditioning signals via Hyperfan-initialized hypernetworks. For MoWM, each HYLA instance generates a different transition matrix:
\begin{equation}
W_m = \mathrm{HYLA}_m(\phi(s_t), \phi(a_t))
\end{equation}

The generated weights replace the fixed $\phi_A$ in the transition equation, making the transformation state-dependent and context-aware.

\subsection{HMM-VSA Rules as World Models}

Each HMM-VSA rule~\cite{grillcheese2025hmm} with transition matrix $A_m$ and emission model $B_m$ \emph{is} a world model: given the current state distribution (forward algorithm posteriors), $A_m$ predicts the next-state distribution. The key insight is that different rules learn different transition patterns---\texttt{distribute\_three} learns periodic dynamics, \texttt{constant} learns equilibria, \texttt{progression} learns monotonic trends. In MoWM, each rule is instantiated as a full world model with its own HYLA-generated transitions.

\section{MoWM Architecture}

\subsection{Multiple Hypernetwork World Models}

Given $M$ HYLA hypernetworks, each produces a state-and-action-dependent transformation:
\begin{equation}
\Delta_m(s, a) = \mathrm{reshape}\!\left(\mathrm{HYLA}_m(\phi(s) \| \phi(a)),\; (k, l)\right)
\end{equation}
where $\|$ denotes concatenation and reshape converts the output to a block-code vector. The predicted next state for world model $m$ is:
\begin{equation}
\hat{s}_{t+1}^{(m)} = \mathrm{bind}(\phi(s_t), \Delta_m(s_t, a_t))
\end{equation}

Each world model preserves the algebraic properties of VSA binding (invertibility, compositionality, magnitude preservation) because $\Delta_m$ is a block-code vector and bind is per-block circular convolution.

\subsection{DSelect-k Gating}

The DSelect-k gate~\cite{hazimeh2021dselect} selects exactly $k$ out of $M$ world models:
\begin{equation}
w_1, \ldots, w_M = \mathrm{DSelect\text{-}k}(\phi(s_t) \| \phi(a_t);\; \theta_\mathrm{gate})
\end{equation}
where exactly $k$ weights are non-zero (enforced by smooth-step activation). The mixture prediction is:
\begin{equation}
\hat{s}_{t+1} = \sum_{m:\, w_m > 0} w_m \cdot \hat{s}_{t+1}^{(m)}
\end{equation}

The entropy regularizer prevents all queries from routing to the same expert, and the diversity loss on HMM transition matrices encourages specialization.

\subsection{Multi-Step Planning via Algebraic Composition}

For planning horizon $H$, we compose transformations:
\begin{equation}
\hat{s}_{t+H} = \phi(s_t) \odot \prod_{j=0}^{H-1} \Delta_{m^*_j}(\hat{s}_{t+j}, a_{t+j})
\end{equation}
where $m^*_j = \arg\max_m w_m(\hat{s}_{t+j}, a_{t+j})$ at each step. In phase space~\cite{chung2025geometric}:
\begin{equation}
\Theta_s^\top \hat{s}_{t+H} = \Theta_s^\top s_t + \sum_{j=0}^{H-1} \mathrm{phase}(\Delta_{m^*_j}) \pmod{2\pi}
\end{equation}

Cleanup is applied every $C$ steps ($C=2$ in our experiments) via \texttt{BlockCodeOps.similarity} against the state codebook, preventing error accumulation.

\subsection{Contrastive Value Estimation}

CVL~\cite{mazoure2022contrastive} estimates Q-values from the mixture's occupancy measure:
\begin{equation}
Q_\mathrm{MoWM}(s, a) = \frac{1}{1-\gamma} \sum_{m:\, w_m > 0} w_m \cdot \mathbb{E}_{\Delta t}\!\left[r(\hat{s}_{t+\Delta t}^{(m)}) \cdot e^{f(s, a, \hat{s}_{t+\Delta t}^{(m)})}\right]
\end{equation}
where $f$ is the contrastive critic learned via InfoNCE. Each world model contributes to the Q-value proportional to its gate weight, and the mixture Q-value naturally accounts for multi-modal dynamics.

\subsection{Training}

The full loss:
\begin{align}
\mathcal{L} &= \lambda_\mathrm{bind}\, \mathcal{L}_\mathrm{bind} + \lambda_\mathrm{inv}\, \mathcal{L}_\mathrm{inv} + \lambda_\mathrm{ortho}\, \mathcal{L}_\mathrm{ortho} \notag \\
&\quad + \lambda_\mathrm{div}\, \mathcal{L}_\mathrm{div} + \lambda_\mathrm{ent}\, \mathcal{L}_\mathrm{ent} + \lambda_\mathrm{CVL}\, \mathcal{L}_\mathrm{CVL}
\end{align}
where:
\begin{itemize}
\item $\mathcal{L}_\mathrm{bind}$: binding loss per world model (each must predict correctly when selected).
\item $\mathcal{L}_\mathrm{inv}$: invertibility of each $\Delta_m$ ($\Delta_m \odot \Delta_m^{-1} \approx \mathrm{identity}$).
\item $\mathcal{L}_\mathrm{ortho}$: orthogonality of state embeddings.
\item $\mathcal{L}_\mathrm{div}$: diversity of HMM transition matrices across world models.
\item $\mathcal{L}_\mathrm{ent}$: DSelect-k entropy regularizer (encourage discrete selection).
\item $\mathcal{L}_\mathrm{CVL}$: contrastive value loss (InfoNCE + partition regularizer).
\end{itemize}

DisARM provides gradients through the discrete DSelect-k selection and block-code discretization. The Surprise-Momentum optimizer modulates learning rate by prediction surprise, and the bandits-based rule explorer determines which world models need more training data.

\subsection{Implementation}

All MoWM components are implemented in the open-source \texttt{grilly} library (Apache 2.0), providing hardware-agnostic Vulkan compute shaders for consumer GPUs (12~GB VRAM target). The \texttt{cubemind} package includes 22 modules and 626 passing tests, all in pure NumPy with optional Vulkan acceleration via Buffer Device Address (BDA) for dynamic GPU memory scaling. Code is available at \url{https://github.com/grillcheese-ai/grilly}.

Robust training extensions (CIW/CICW for noisy labels, DROPS for class imbalance, Staircase DP for privacy, R2B for fairness) are available in the \texttt{cubemind} package but are orthogonal to the core architecture and evaluated separately.

\section{Experiments}

\subsection{Component Validation}

We validated each MoWM component individually before full environment integration. All benchmarks were run on consumer GPU hardware (12~GB VRAM).

\textbf{Block-code operations.} Block-code binding executed at 57~$\mu$s/op and similarity search at 4~$\mu$s/op. These are the core operations in each world model's transition function---the combined 61~$\mu$s/op per bind-then-cleanup cycle confirms that MoWM's algebraic planning (Section~III-C) is feasible at interactive rates even for long rollout horizons.

\textbf{State processing: Combiner-Axial vs FNet mixing.} FNet mixing was 2$\times$ faster than Combiner-Axial attention across sequence lengths:

\begin{table}[t]
\centering
\caption{Mixing method latency comparison.}
\label{tab:mixing}
\begin{tabular}{lcc}
\toprule
Mixing Method & $L=256$ & $L=1024$ \\
\midrule
Combiner-Axial & 0.84 ms & 2.13 ms \\
FNet & 0.31 ms & 1.20 ms \\
\bottomrule
\end{tabular}
\end{table}

FNet's speed advantage makes it the default state-processing layer in MoWM, where the mixing step runs $M$ times per forward pass (once per active world model). At $M=4$, $k=2$, FNet saves approximately 1.0~ms per step at $L=256$ compared to Combiner-Axial.

\textbf{HYLA hypernetwork stability.} We measured weight variance propagation through 5-layer HYLA stacks that generate each world model's transition weights. Hyperfan initialization without MIP achieved variance $7.9 \times 10^{-7}$ after 5 layers---effectively zero drift, confirming that deep hypernetwork chains remain stable. Xavier initialization with MIP was also stable (variance $2.1 \times 10^{-2}$ after 5 layers), providing a usable fallback. These results validate that $M$ independent HYLA hypernetworks can each generate transition matrices without variance explosion or collapse.

\textbf{DisARM gradient estimation.} Gradient variance through discrete DSelect-k selection was measured across three estimators:

\begin{table}[t]
\centering
\caption{Gradient variance comparison across estimators.}
\label{tab:disarm}
\begin{tabular}{lc}
\toprule
Estimator & Variance \\
\midrule
DisARM & 0.234 \\
ARM & 0.279 \\
REINFORCE & 1.751 \\
\bottomrule
\end{tabular}
\end{table}

DisARM achieved 6.4$\times$ lower variance than REINFORCE and 1.2$\times$ lower than ARM, confirming it as the correct gradient estimator for training through the discrete model selection in DSelect-k gating.

\textbf{VS-Graph perception frontend.} VS-Graph achieved 100\% classification accuracy on star-vs-chain topology discrimination tasks, validating it as the perception frontend for encoding graph-structured inputs into block-code vectors for the MoWM pipeline (Section~V).

\subsection{Grid World (Direct Comparison with Chung et al.)}

\begin{table}[t]
\centering
\caption{Grid world results comparing MLP-Large and FHRR baselines.}
\label{tab:gridworld}
\begin{tabular}{lcccc}
\toprule
Model & 1-step Acc & Zero-shot & Roll-20 & Roll-20+Clean \\
\midrule
MLP-Large & 80.25\% & 1.25\% & 6.2\% & 8.4\% \\
FHRR (Chung) & 96.3\% & 87.5\% & 34.6\% & 61.4\% \\
\bottomrule
\end{tabular}
\end{table}

Full grid-world evaluation of MoWM ($M=4$, $k=2$) and MoWM ($M=8$, $k=2$) configurations is in progress. The component results above confirm that the computational budget per transition step (block-code bind at 57~$\mu$s + FNet mixing at 0.31~ms) is well within the latency envelope required for interactive rollout planning.

\subsection{Multi-Modal Environments}

Environments where the same action has different effects depending on context:
\begin{itemize}
\item Grid world with locked/unlocked doors.
\item Grid world with varying terrain (ice, mud, normal).
\item Simple physics simulation (gravity changes per region).
\end{itemize}

Multi-modal environments are the primary target for MoWM's advantage over single-model baselines---different world models specialize in different terrain types via the diversity loss on HMM transition matrices.

\subsection{Ablation Studies}

\begin{table*}[t]
\centering
\caption{Ablation study of MoWM components.}
\label{tab:ablation}
\begin{tabular}{ll}
\toprule
Ablation & Observed / Measured Effect \\
\midrule
$M=1$ (single world model) & Reduces to Chung et al.\ baseline \\
$k=M$ (all models active) & No sparsity, slower inference, possible mode collapse \\
No DSelect-k (soft weights) & All models contribute, less specialization \\
No cleanup & Error accumulation on long rollouts \\
No diversity loss & Mode collapse---all world models converge to same transitions \\
No CVL (TD-learning instead) & Less stable Q-values, requires bootstrapping \\
Block codes vs FHRR & Block codes: 57~$\mu$s/op bind, 4~$\mu$s/op similarity, no magnitude drift; FHRR: continuous but drifts \\
Combiner vs FNet mixing & FNet 2$\times$ faster (0.31~ms vs 0.84~ms at $L{=}256$; 1.20~ms vs 2.13~ms at $L{=}1024$) \\
DisARM vs ARM vs REINFORCE & Variance: DisARM 0.234 $<$ ARM 0.279 $<$ REINFORCE 1.751 \\
HYLA init: Hyperfan+noMIP & Stable: var $= 7.9 \times 10^{-7}$ after 5 layers \\
HYLA init: Xavier+MIP & Stable: var $= 2.1 \times 10^{-2}$ after 5 layers \\
\bottomrule
\end{tabular}
\end{table*}

\subsection{Specialization Analysis}

We analyze what each world model learns through:
\begin{itemize}
\item Transition matrix heatmaps per model.
\item t-SNE of state embeddings colored by which model activates.
\item Gate activation patterns across state space.
\item Information-flow analysis (fact association probes) through each world model.
\end{itemize}

\section{Related Work}

\textbf{Single-model VSA world models.} Chung et al.~\cite{chung2025geometric} demonstrated VSA binding as transition with FHRR encoders. We extend to multiple hypernetwork-generated models.

\textbf{VSA-based graph perception.} Poursiami et al.~\cite{poursiami2025vsgraph} proposed VS-Graph, achieving GNN-competitive graph classification 450$\times$ faster via Spike Diffusion (topology-aware node ranking) and Associative Message Passing (idempotent OR aggregation). We integrated VS-Graph as the perception frontend for MoWM: graph-structured inputs are encoded into block-code vectors via Spike Diffusion + block-code bundling, feeding directly into the hypernetwork world models. In our validation, VS-Graph achieved 100\% classification accuracy on star-vs-chain topology discrimination, confirming the fully VSA end-to-end pipeline.

\textbf{Mixture of Experts.} Shazeer et al.~\cite{shazeer2017moe} introduced sparsely-gated MoE for language models. Hazimeh et al.~\cite{hazimeh2021dselect} proposed DSelect-k for differentiable top-$k$ selection. We apply MoE to world models, where each expert is a complete transition function.

\textbf{Multi-model RL.} Chua et al.~\cite{chua2018pets} used an ensemble of dynamics models for model-based RL (PETS), but without algebraic structure or compositional generalization. Lee et al.~\cite{lee2020sunrise} used network ensembles for uncertainty estimation. Our approach differs in that each model has VSA structure enabling composition and cleanup.

\textbf{HMM-VSA.} Grillcheese Research~\cite{grillcheese2025hmm} proposed learnable HMM rules grounded in GFSA theory. MoWM instantiates each rule as a full hypernetwork-generated world model.

\section{Conclusion}

MoWM demonstrates that combining multiple algebraically-structured world models via sparse expert selection produces robust, interpretable, and generalizable environment dynamics. Each HYLA hypernetwork generates a different VSA transition function, the DSelect-k gate learns when to apply each, and the block-code algebra enables compositional planning with cleanup-based error correction. The framework is theoretically grounded in GFSA theory and practically implementable on consumer hardware.

Component validation confirmed the feasibility of every core operation: block-code binding at 57~$\mu$s/op with 4~$\mu$s/op similarity search, FNet state mixing at 2$\times$ the speed of Combiner-Axial attention, HYLA hypernetworks stable through 5 layers (variance $7.9 \times 10^{-7}$ with Hyperfan+noMIP), DisARM gradient variance 6.4$\times$ lower than REINFORCE for training through discrete model selection, and VS-Graph achieving 100\% accuracy on topology classification as the perception frontend. These results establish that the MoWM pipeline operates well within interactive latency budgets on consumer hardware.

The path from single-model world modeling~\cite{chung2025geometric} to MoWM parallels the evolution from single attention heads to multi-head attention in transformers: multiple specialized views of the dynamics, combined via learned gating, produce richer representations than any single view alone.

\begin{thebibliography}{99}

\bibitem{andrychowicz2017her}
M.~Andrychowicz et~al., ``Hindsight experience replay,'' in \emph{Proc. Advances in Neural Information Processing Systems (NeurIPS)}, 2017.

\bibitem{chung2025geometric}
W.~Y. Chung, C.~Yeung, H.~J. Lillemark, Z.~Zou, X.~Liu, and M.~Imani, ``Geometric priors for generalizable world models via vector symbolic architecture,'' in \emph{NeurReps, Under Review}, 2025.

\bibitem{chua2018pets}
K.~Chua, R.~Calandra, R.~McAllister, and S.~Levine, ``Deep reinforcement learning in a handful of trials using probabilistic dynamics models,'' in \emph{Proc. Advances in Neural Information Processing Systems (NeurIPS)}, 2018.

\bibitem{grillcheese2025hmm}
N.~Cloutier, ``Learnable reasoning rules for neuro-vector-symbolic architectures via hidden Markov models,'' \emph{Zenodo}, 2025. [Online]. Available: \url{https://doi.org/10.5281/zenodo.19190708}

\bibitem{ha2018world}
D.~Ha and J.~Schmidhuber, ``World models,'' \emph{arXiv preprint arXiv:1803.10122}, 2018.

\bibitem{hafner2020mastering}
D.~Hafner, T.~Lillicrap, M.~Norouzi, and J.~Ba, ``Mastering Atari with discrete world models,'' \emph{arXiv preprint arXiv:2010.02193}, 2020.

\bibitem{hansen2023tdmpc2}
N.~Hansen, H.~Su, and X.~Wang, ``TD-MPC2: Scalable, robust world models for continuous control,'' \emph{arXiv preprint arXiv:2310.16828}, 2023.

\bibitem{hazimeh2021dselect}
H.~Hazimeh et~al., ``DSelect-k: Differentiable selection in the mixture of experts,'' in \emph{Proc. Advances in Neural Information Processing Systems (NeurIPS)}, 2021.

\bibitem{hersche2023nvsa}
M.~Hersche et~al., ``A neuro-vector-symbolic architecture for solving Raven's progressive matrices,'' \emph{Nature Machine Intelligence}, 2023.

\bibitem{johnson2020learning}
D.~D. Johnson, H.~Larochelle, and D.~Tarlow, ``Learning graph structure with a finite-state automaton layer,'' in \emph{Proc. Advances in Neural Information Processing Systems (NeurIPS)}, 2020.

\bibitem{lee2020sunrise}
K.~Lee et~al., ``SUNRISE: A simple unified framework for ensemble learning in deep reinforcement learning,'' in \emph{Proc. International Conference on Machine Learning (ICML)}, 2020.

\bibitem{mazoure2022contrastive}
B.~Mazoure, B.~Eysenbach, O.~Nachum, and J.~Tompson, ``Contrastive value learning,'' in \emph{NeurIPS Workshop}, 2022.

\bibitem{poursiami2025vsgraph}
A.~Poursiami et~al., ``VS-Graph,'' 2025.

\bibitem{shazeer2017moe}
N.~Shazeer et~al., ``Outrageously large neural networks: The sparsely-gated mixture-of-experts layer,'' in \emph{Proc. International Conference on Learning Representations (ICLR)}, 2017.

\end{thebibliography}

\end{document}
